\textbf{A continuación se presenta el detalle departamental del crecimiento de la productividad laboral.} Para cada departamento se reportan el PIB real, el número de ocupados, el PIB por trabajador, las horas semanales por trabajador y el PIB por hora trabajada al inicio y al final del periodo.

\subsection{Quindío}
\textbf{Entre 2010 y 2024pr, el PIB por hora trabajada de Quindío creció 4,63\% anual.} El PIB por trabajador creció 3,78\% anual, mientras que las horas semanales por trabajador cambiaron -0,81\% anual. Esta diferencia muestra si la productividad medida por trabajador se mueve en línea con la productividad por hora o si está afectada por cambios en la intensidad horaria.
\begin{table}[H]
\centering
\caption{Quindío: PIB, ocupados, horas y productividad laboral, 2010--2024pr}
\label{tab:departamento_quindio_productividad}
\scriptsize
\begin{tabular}{lrrr}
\toprule
Indicador & 2010 & 2024pr & Crec. anual \\
\midrule
PIB real (Billones de pesos de 2015) & 5,1 & 8,0 & 3,24\% \\
Ocupados (Millones) & 0,25 & 0,24 & -0,53\% \\
PIB por trabajador (Millones de pesos de 2015) & 20,1 & 33,8 & 3,78\% \\
Horas semanales por trabajador & 48,3 & 43,1 & -0,81\% \\
PIB por hora trabajada (Miles de pesos de 2015) & 8,0 & 15,1 & 4,63\% \\
\bottomrule
\end{tabular}
\end{table}

\subsection{Caquetá}
\textbf{Entre 2010 y 2024pr, el PIB por hora trabajada de Caquetá creció 4,17\% anual.} El PIB por trabajador creció 3,71\% anual, mientras que las horas semanales por trabajador cambiaron -0,45\% anual. Esta diferencia muestra si la productividad medida por trabajador se mueve en línea con la productividad por hora o si está afectada por cambios en la intensidad horaria.
\begin{table}[H]
\centering
\caption{Caquetá: PIB, ocupados, horas y productividad laboral, 2010--2024pr}
\label{tab:departamento_caqueta_productividad}
\scriptsize
\begin{tabular}{lrrr}
\toprule
Indicador & 2010 & 2024pr & Crec. anual \\
\midrule
PIB real (Billones de pesos de 2015) & 2,6 & 3,9 & 3,01\% \\
Ocupados (Millones) & 0,17 & 0,16 & -0,67\% \\
PIB por trabajador (Millones de pesos de 2015) & 14,9 & 24,9 & 3,71\% \\
Horas semanales por trabajador & 48,7 & 45,7 & -0,45\% \\
PIB por hora trabajada (Miles de pesos de 2015) & 5,9 & 10,5 & 4,17\% \\
\bottomrule
\end{tabular}
\end{table}

\subsection{Risaralda}
\textbf{Entre 2010 y 2024pr, el PIB por hora trabajada de Risaralda creció 3,51\% anual.} El PIB por trabajador creció 3,14\% anual, mientras que las horas semanales por trabajador cambiaron -0,36\% anual. Esta diferencia muestra si la productividad medida por trabajador se mueve en línea con la productividad por hora o si está afectada por cambios en la intensidad horaria.
\begin{table}[H]
\centering
\caption{Risaralda: PIB, ocupados, horas y productividad laboral, 2010--2024pr}
\label{tab:departamento_risaralda_productividad}
\scriptsize
\begin{tabular}{lrrr}
\toprule
Indicador & 2010 & 2024pr & Crec. anual \\
\midrule
PIB real (Billones de pesos de 2015) & 10,5 & 16,3 & 3,21\% \\
Ocupados (Millones) & 0,42 & 0,42 & 0,07\% \\
PIB por trabajador (Millones de pesos de 2015) & 25,1 & 38,7 & 3,14\% \\
Horas semanales por trabajador & 47,7 & 45,4 & -0,36\% \\
PIB por hora trabajada (Miles de pesos de 2015) & 10,1 & 16,4 & 3,51\% \\
\bottomrule
\end{tabular}
\end{table}

\subsection{Chocó}
\textbf{Entre 2010 y 2024pr, el PIB por hora trabajada de Chocó creció 3,20\% anual.} El PIB por trabajador creció 2,95\% anual, mientras que las horas semanales por trabajador cambiaron -0,24\% anual. Esta diferencia muestra si la productividad medida por trabajador se mueve en línea con la productividad por hora o si está afectada por cambios en la intensidad horaria.
\begin{table}[H]
\centering
\caption{Chocó: PIB, ocupados, horas y productividad laboral, 2010--2024pr}
\label{tab:departamento_choco_productividad}
\scriptsize
\begin{tabular}{lrrr}
\toprule
Indicador & 2010 & 2024pr & Crec. anual \\
\midrule
PIB real (Billones de pesos de 2015) & 3,8 & 3,8 & -0,05\% \\
Ocupados (Millones) & 0,20 & 0,13 & -2,91\% \\
PIB por trabajador (Millones de pesos de 2015) & 19,0 & 28,6 & 2,95\% \\
Horas semanales por trabajador & 47,7 & 46,1 & -0,24\% \\
PIB por hora trabajada (Miles de pesos de 2015) & 7,7 & 11,9 & 3,20\% \\
\bottomrule
\end{tabular}
\end{table}

\subsection{Caldas}
\textbf{Entre 2010 y 2024pr, el PIB por hora trabajada de Caldas creció 2,57\% anual.} El PIB por trabajador creció 2,14\% anual, mientras que las horas semanales por trabajador cambiaron -0,41\% anual. Esta diferencia muestra si la productividad medida por trabajador se mueve en línea con la productividad por hora o si está afectada por cambios en la intensidad horaria.
\begin{table}[H]
\centering
\caption{Caldas: PIB, ocupados, horas y productividad laboral, 2010--2024pr}
\label{tab:departamento_caldas_productividad}
\scriptsize
\begin{tabular}{lrrr}
\toprule
Indicador & 2010 & 2024pr & Crec. anual \\
\midrule
PIB real (Billones de pesos de 2015) & 10,7 & 15,7 & 2,82\% \\
Ocupados (Millones) & 0,39 & 0,43 & 0,66\% \\
PIB por trabajador (Millones de pesos de 2015) & 27,2 & 36,6 & 2,14\% \\
Horas semanales por trabajador & 47,9 & 45,3 & -0,41\% \\
PIB por hora trabajada (Miles de pesos de 2015) & 10,9 & 15,6 & 2,57\% \\
\bottomrule
\end{tabular}
\end{table}

\subsection{Valle del Cauca}
\textbf{Entre 2010 y 2024pr, el PIB por hora trabajada de Valle del Cauca creció 2,50\% anual.} El PIB por trabajador creció 2,15\% anual, mientras que las horas semanales por trabajador cambiaron -0,35\% anual. Esta diferencia muestra si la productividad medida por trabajador se mueve en línea con la productividad por hora o si está afectada por cambios en la intensidad horaria.
\begin{table}[H]
\centering
\caption{Valle del Cauca: PIB, ocupados, horas y productividad laboral, 2010--2024pr}
\label{tab:departamento_valle_del_cauca_productividad}
\scriptsize
\begin{tabular}{lrrr}
\toprule
Indicador & 2010 & 2024pr & Crec. anual \\
\midrule
PIB real (Billones de pesos de 2015) & 63,8 & 99,5 & 3,23\% \\
Ocupados (Millones) & 1,58 & 1,83 & 1,05\% \\
PIB por trabajador (Millones de pesos de 2015) & 40,5 & 54,5 & 2,15\% \\
Horas semanales por trabajador & 46,9 & 44,7 & -0,35\% \\
PIB por hora trabajada (Miles de pesos de 2015) & 16,6 & 23,4 & 2,50\% \\
\bottomrule
\end{tabular}
\end{table}

\subsection{Sucre}
\textbf{Entre 2010 y 2024pr, el PIB por hora trabajada de Sucre creció 2,45\% anual.} El PIB por trabajador creció 1,67\% anual, mientras que las horas semanales por trabajador cambiaron -0,77\% anual. Esta diferencia muestra si la productividad medida por trabajador se mueve en línea con la productividad por hora o si está afectada por cambios en la intensidad horaria.
\begin{table}[H]
\centering
\caption{Sucre: PIB, ocupados, horas y productividad laboral, 2010--2024pr}
\label{tab:departamento_sucre_productividad}
\scriptsize
\begin{tabular}{lrrr}
\toprule
Indicador & 2010 & 2024pr & Crec. anual \\
\midrule
PIB real (Billones de pesos de 2015) & 5,2 & 8,1 & 3,23\% \\
Ocupados (Millones) & 0,30 & 0,37 & 1,54\% \\
PIB por trabajador (Millones de pesos de 2015) & 17,5 & 22,0 & 1,67\% \\
Horas semanales por trabajador & 45,4 & 40,8 & -0,77\% \\
PIB por hora trabajada (Miles de pesos de 2015) & 7,4 & 10,4 & 2,45\% \\
\bottomrule
\end{tabular}
\end{table}

\subsection{Bogotá D.C.}
\textbf{Entre 2010 y 2024pr, el PIB por hora trabajada de Bogotá D.C. creció 2,44\% anual.} El PIB por trabajador creció 1,98\% anual, mientras que las horas semanales por trabajador cambiaron -0,45\% anual. Esta diferencia muestra si la productividad medida por trabajador se mueve en línea con la productividad por hora o si está afectada por cambios en la intensidad horaria.
\begin{table}[H]
\centering
\caption{Bogotá D.C.: PIB, ocupados, horas y productividad laboral, 2010--2024pr}
\label{tab:departamento_bogota_dc_productividad}
\scriptsize
\begin{tabular}{lrrr}
\toprule
Indicador & 2010 & 2024pr & Crec. anual \\
\midrule
PIB real (Billones de pesos de 2015) & 167,4 & 267,1 & 3,39\% \\
Ocupados (Millones) & 3,27 & 3,96 & 1,39\% \\
PIB por trabajador (Millones de pesos de 2015) & 51,2 & 67,4 & 1,98\% \\
Horas semanales por trabajador & 48,3 & 45,4 & -0,45\% \\
PIB por hora trabajada (Miles de pesos de 2015) & 20,4 & 28,6 & 2,44\% \\
\bottomrule
\end{tabular}
\end{table}

\subsection{Bolívar}
\textbf{Entre 2010 y 2024pr, el PIB por hora trabajada de Bolívar creció 2,34\% anual.} El PIB por trabajador creció 1,69\% anual, mientras que las horas semanales por trabajador cambiaron -0,64\% anual. Esta diferencia muestra si la productividad medida por trabajador se mueve en línea con la productividad por hora o si está afectada por cambios en la intensidad horaria.
\begin{table}[H]
\centering
\caption{Bolívar: PIB, ocupados, horas y productividad laboral, 2010--2024pr}
\label{tab:departamento_bolivar_productividad}
\scriptsize
\begin{tabular}{lrrr}
\toprule
Indicador & 2010 & 2024pr & Crec. anual \\
\midrule
PIB real (Billones de pesos de 2015) & 23,1 & 35,9 & 3,21\% \\
Ocupados (Millones) & 0,65 & 0,80 & 1,50\% \\
PIB por trabajador (Millones de pesos de 2015) & 35,7 & 45,1 & 1,69\% \\
Horas semanales por trabajador & 48,0 & 43,9 & -0,64\% \\
PIB por hora trabajada (Miles de pesos de 2015) & 14,3 & 19,8 & 2,34\% \\
\bottomrule
\end{tabular}
\end{table}

\subsection{Tolima}
\textbf{Entre 2010 y 2024pr, el PIB por hora trabajada de Tolima creció 2,17\% anual.} El PIB por trabajador creció 2,04\% anual, mientras que las horas semanales por trabajador cambiaron -0,12\% anual. Esta diferencia muestra si la productividad medida por trabajador se mueve en línea con la productividad por hora o si está afectada por cambios en la intensidad horaria.
\begin{table}[H]
\centering
\caption{Tolima: PIB, ocupados, horas y productividad laboral, 2010--2024pr}
\label{tab:departamento_tolima_productividad}
\scriptsize
\begin{tabular}{lrrr}
\toprule
Indicador & 2010 & 2024pr & Crec. anual \\
\midrule
PIB real (Billones de pesos de 2015) & 14,8 & 20,3 & 2,28\% \\
Ocupados (Millones) & 0,49 & 0,51 & 0,23\% \\
PIB por trabajador (Millones de pesos de 2015) & 30,2 & 40,0 & 2,04\% \\
Horas semanales por trabajador & 46,3 & 45,5 & -0,12\% \\
PIB por hora trabajada (Miles de pesos de 2015) & 12,5 & 16,9 & 2,17\% \\
\bottomrule
\end{tabular}
\end{table}

\subsection{Cauca}
\textbf{Entre 2010 y 2024pr, el PIB por hora trabajada de Cauca creció 2,05\% anual.} El PIB por trabajador creció 1,44\% anual, mientras que las horas semanales por trabajador cambiaron -0,60\% anual. Esta diferencia muestra si la productividad medida por trabajador se mueve en línea con la productividad por hora o si está afectada por cambios en la intensidad horaria.
\begin{table}[H]
\centering
\caption{Cauca: PIB, ocupados, horas y productividad laboral, 2010--2024pr}
\label{tab:departamento_cauca_productividad}
\scriptsize
\begin{tabular}{lrrr}
\toprule
Indicador & 2010 & 2024pr & Crec. anual \\
\midrule
PIB real (Billones de pesos de 2015) & 10,5 & 17,4 & 3,66\% \\
Ocupados (Millones) & 0,48 & 0,65 & 2,19\% \\
PIB por trabajador (Millones de pesos de 2015) & 21,8 & 26,6 & 1,44\% \\
Horas semanales por trabajador & 40,7 & 37,4 & -0,60\% \\
PIB por hora trabajada (Miles de pesos de 2015) & 10,3 & 13,7 & 2,05\% \\
\bottomrule
\end{tabular}
\end{table}

\subsection{Santander}
\textbf{Entre 2010 y 2024pr, el PIB por hora trabajada de Santander creció 1,96\% anual.} El PIB por trabajador creció 1,71\% anual, mientras que las horas semanales por trabajador cambiaron -0,24\% anual. Esta diferencia muestra si la productividad medida por trabajador se mueve en línea con la productividad por hora o si está afectada por cambios en la intensidad horaria.
\begin{table}[H]
\centering
\caption{Santander: PIB, ocupados, horas y productividad laboral, 2010--2024pr}
\label{tab:departamento_santander_productividad}
\scriptsize
\begin{tabular}{lrrr}
\toprule
Indicador & 2010 & 2024pr & Crec. anual \\
\midrule
PIB real (Billones de pesos de 2015) & 42,8 & 62,6 & 2,75\% \\
Ocupados (Millones) & 0,88 & 1,01 & 1,02\% \\
PIB por trabajador (Millones de pesos de 2015) & 48,8 & 62,0 & 1,71\% \\
Horas semanales por trabajador & 47,2 & 45,6 & -0,24\% \\
PIB por hora trabajada (Miles de pesos de 2015) & 19,9 & 26,1 & 1,96\% \\
\bottomrule
\end{tabular}
\end{table}

\subsection{Nariño}
\textbf{Entre 2010 y 2024pr, el PIB por hora trabajada de Nariño creció 1,92\% anual.} El PIB por trabajador creció 0,97\% anual, mientras que las horas semanales por trabajador cambiaron -0,93\% anual. Esta diferencia muestra si la productividad medida por trabajador se mueve en línea con la productividad por hora o si está afectada por cambios en la intensidad horaria.
\begin{table}[H]
\centering
\caption{Nariño: PIB, ocupados, horas y productividad laboral, 2010--2024pr}
\label{tab:departamento_narino_productividad}
\scriptsize
\begin{tabular}{lrrr}
\toprule
Indicador & 2010 & 2024pr & Crec. anual \\
\midrule
PIB real (Billones de pesos de 2015) & 9,3 & 14,7 & 3,34\% \\
Ocupados (Millones) & 0,60 & 0,82 & 2,34\% \\
PIB por trabajador (Millones de pesos de 2015) & 15,5 & 17,8 & 0,97\% \\
Horas semanales por trabajador & 42,7 & 37,5 & -0,93\% \\
PIB por hora trabajada (Miles de pesos de 2015) & 7,0 & 9,1 & 1,92\% \\
\bottomrule
\end{tabular}
\end{table}

\subsection{Córdoba}
\textbf{Entre 2010 y 2024pr, el PIB por hora trabajada de Córdoba creció 1,71\% anual.} El PIB por trabajador creció 0,90\% anual, mientras que las horas semanales por trabajador cambiaron -0,80\% anual. Esta diferencia muestra si la productividad medida por trabajador se mueve en línea con la productividad por hora o si está afectada por cambios en la intensidad horaria.
\begin{table}[H]
\centering
\caption{Córdoba: PIB, ocupados, horas y productividad laboral, 2010--2024pr}
\label{tab:departamento_cordoba_productividad}
\scriptsize
\begin{tabular}{lrrr}
\toprule
Indicador & 2010 & 2024pr & Crec. anual \\
\midrule
PIB real (Billones de pesos de 2015) & 11,6 & 17,0 & 2,78\% \\
Ocupados (Millones) & 0,60 & 0,77 & 1,86\% \\
PIB por trabajador (Millones de pesos de 2015) & 19,4 & 22,0 & 0,90\% \\
Horas semanales por trabajador & 44,4 & 39,7 & -0,80\% \\
PIB por hora trabajada (Miles de pesos de 2015) & 8,4 & 10,7 & 1,71\% \\
\bottomrule
\end{tabular}
\end{table}

\subsection{Antioquia}
\textbf{Entre 2010 y 2024pr, el PIB por hora trabajada de Antioquia creció 1,58\% anual.} El PIB por trabajador creció 1,27\% anual, mientras que las horas semanales por trabajador cambiaron -0,30\% anual. Esta diferencia muestra si la productividad medida por trabajador se mueve en línea con la productividad por hora o si está afectada por cambios en la intensidad horaria.
\begin{table}[H]
\centering
\caption{Antioquia: PIB, ocupados, horas y productividad laboral, 2010--2024pr}
\label{tab:departamento_antioquia_productividad}
\scriptsize
\begin{tabular}{lrrr}
\toprule
Indicador & 2010 & 2024pr & Crec. anual \\
\midrule
PIB real (Billones de pesos de 2015) & 88,5 & 149,2 & 3,80\% \\
Ocupados (Millones) & 2,22 & 3,13 & 2,50\% \\
PIB por trabajador (Millones de pesos de 2015) & 39,9 & 47,7 & 1,27\% \\
Horas semanales por trabajador & 47,0 & 45,1 & -0,30\% \\
PIB por hora trabajada (Miles de pesos de 2015) & 16,3 & 20,3 & 1,58\% \\
\bottomrule
\end{tabular}
\end{table}

\subsection{Boyacá}
\textbf{Entre 2010 y 2024pr, el PIB por hora trabajada de Boyacá creció 1,43\% anual.} El PIB por trabajador creció 1,16\% anual, mientras que las horas semanales por trabajador cambiaron -0,27\% anual. Esta diferencia muestra si la productividad medida por trabajador se mueve en línea con la productividad por hora o si está afectada por cambios en la intensidad horaria.
\begin{table}[H]
\centering
\caption{Boyacá: PIB, ocupados, horas y productividad laboral, 2010--2024pr}
\label{tab:departamento_boyaca_productividad}
\scriptsize
\begin{tabular}{lrrr}
\toprule
Indicador & 2010 & 2024pr & Crec. anual \\
\midrule
PIB real (Billones de pesos de 2015) & 17,3 & 25,2 & 2,70\% \\
Ocupados (Millones) & 0,43 & 0,53 & 1,52\% \\
PIB por trabajador (Millones de pesos de 2015) & 40,4 & 47,5 & 1,16\% \\
Horas semanales por trabajador & 42,0 & 40,4 & -0,27\% \\
PIB por hora trabajada (Miles de pesos de 2015) & 18,5 & 22,6 & 1,43\% \\
\bottomrule
\end{tabular}
\end{table}

\subsection{Norte de Santander}
\textbf{Entre 2010 y 2024pr, el PIB por hora trabajada de Norte de Santander creció 1,41\% anual.} El PIB por trabajador creció 1,23\% anual, mientras que las horas semanales por trabajador cambiaron -0,17\% anual. Esta diferencia muestra si la productividad medida por trabajador se mueve en línea con la productividad por hora o si está afectada por cambios en la intensidad horaria.
\begin{table}[H]
\centering
\caption{Norte de Santander: PIB, ocupados, horas y productividad laboral, 2010--2024pr}
\label{tab:departamento_norte_de_santander_productividad}
\scriptsize
\begin{tabular}{lrrr}
\toprule
Indicador & 2010 & 2024pr & Crec. anual \\
\midrule
PIB real (Billones de pesos de 2015) & 10,7 & 15,6 & 2,74\% \\
Ocupados (Millones) & 0,52 & 0,64 & 1,49\% \\
PIB por trabajador (Millones de pesos de 2015) & 20,5 & 24,3 & 1,23\% \\
Horas semanales por trabajador & 47,8 & 46,7 & -0,17\% \\
PIB por hora trabajada (Miles de pesos de 2015) & 8,2 & 10,0 & 1,41\% \\
\bottomrule
\end{tabular}
\end{table}

\subsection{Magdalena}
\textbf{Entre 2010 y 2024pr, el PIB por hora trabajada de Magdalena creció 0,88\% anual.} El PIB por trabajador creció 0,16\% anual, mientras que las horas semanales por trabajador cambiaron -0,71\% anual. Esta diferencia muestra si la productividad medida por trabajador se mueve en línea con la productividad por hora o si está afectada por cambios en la intensidad horaria.
\begin{table}[H]
\centering
\caption{Magdalena: PIB, ocupados, horas y productividad laboral, 2010--2024pr}
\label{tab:departamento_magdalena_productividad}
\scriptsize
\begin{tabular}{lrrr}
\toprule
Indicador & 2010 & 2024pr & Crec. anual \\
\midrule
PIB real (Billones de pesos de 2015) & 9,1 & 13,1 & 2,62\% \\
Ocupados (Millones) & 0,39 & 0,55 & 2,46\% \\
PIB por trabajador (Millones de pesos de 2015) & 23,3 & 23,9 & 0,16\% \\
Horas semanales por trabajador & 49,5 & 44,8 & -0,71\% \\
PIB por hora trabajada (Miles de pesos de 2015) & 9,1 & 10,3 & 0,88\% \\
\bottomrule
\end{tabular}
\end{table}

\subsection{Atlántico}
\textbf{Entre 2010 y 2024pr, el PIB por hora trabajada de Atlántico creció 0,85\% anual.} El PIB por trabajador creció 0,46\% anual, mientras que las horas semanales por trabajador cambiaron -0,38\% anual. Esta diferencia muestra si la productividad medida por trabajador se mueve en línea con la productividad por hora o si está afectada por cambios en la intensidad horaria.
\begin{table}[H]
\centering
\caption{Atlántico: PIB, ocupados, horas y productividad laboral, 2010--2024pr}
\label{tab:departamento_atlantico_productividad}
\scriptsize
\begin{tabular}{lrrr}
\toprule
Indicador & 2010 & 2024pr & Crec. anual \\
\midrule
PIB real (Billones de pesos de 2015) & 26,6 & 45,1 & 3,84\% \\
Ocupados (Millones) & 0,76 & 1,21 & 3,36\% \\
PIB por trabajador (Millones de pesos de 2015) & 35,0 & 37,3 & 0,46\% \\
Horas semanales por trabajador & 47,5 & 45,0 & -0,38\% \\
PIB por hora trabajada (Miles de pesos de 2015) & 14,2 & 16,0 & 0,85\% \\
\bottomrule
\end{tabular}
\end{table}

\subsection{Cesar}
\textbf{Entre 2010 y 2024pr, el PIB por hora trabajada de Cesar creció 0,57\% anual.} El PIB por trabajador creció -0,03\% anual, mientras que las horas semanales por trabajador cambiaron -0,60\% anual. Esta diferencia muestra si la productividad medida por trabajador se mueve en línea con la productividad por hora o si está afectada por cambios en la intensidad horaria.
\begin{table}[H]
\centering
\caption{Cesar: PIB, ocupados, horas y productividad laboral, 2010--2024pr}
\label{tab:departamento_cesar_productividad}
\scriptsize
\begin{tabular}{lrrr}
\toprule
Indicador & 2010 & 2024pr & Crec. anual \\
\midrule
PIB real (Billones de pesos de 2015) & 11,6 & 15,2 & 1,97\% \\
Ocupados (Millones) & 0,35 & 0,47 & 2,00\% \\
PIB por trabajador (Millones de pesos de 2015) & 32,7 & 32,6 & -0,03\% \\
Horas semanales por trabajador & 49,2 & 45,2 & -0,60\% \\
PIB por hora trabajada (Miles de pesos de 2015) & 12,8 & 13,8 & 0,57\% \\
\bottomrule
\end{tabular}
\end{table}

\subsection{Meta}
\textbf{Entre 2010 y 2024pr, el PIB por hora trabajada de Meta creció 0,55\% anual.} El PIB por trabajador creció 0,05\% anual, mientras que las horas semanales por trabajador cambiaron -0,50\% anual. Esta diferencia muestra si la productividad medida por trabajador se mueve en línea con la productividad por hora o si está afectada por cambios en la intensidad horaria.
\begin{table}[H]
\centering
\caption{Meta: PIB, ocupados, horas y productividad laboral, 2010--2024pr}
\label{tab:departamento_meta_productividad}
\scriptsize
\begin{tabular}{lrrr}
\toprule
Indicador & 2010 & 2024pr & Crec. anual \\
\midrule
PIB real (Billones de pesos de 2015) & 21,3 & 32,1 & 2,98\% \\
Ocupados (Millones) & 0,34 & 0,51 & 2,93\% \\
PIB por trabajador (Millones de pesos de 2015) & 62,7 & 63,1 & 0,05\% \\
Horas semanales por trabajador & 50,3 & 46,9 & -0,50\% \\
PIB por hora trabajada (Miles de pesos de 2015) & 24,0 & 25,9 & 0,55\% \\
\bottomrule
\end{tabular}
\end{table}

\subsection{Huila}
\textbf{Entre 2010 y 2024pr, el PIB por hora trabajada de Huila creció 0,11\% anual.} El PIB por trabajador creció -0,13\% anual, mientras que las horas semanales por trabajador cambiaron -0,24\% anual. Esta diferencia muestra si la productividad medida por trabajador se mueve en línea con la productividad por hora o si está afectada por cambios en la intensidad horaria.
\begin{table}[H]
\centering
\caption{Huila: PIB, ocupados, horas y productividad laboral, 2010--2024pr}
\label{tab:departamento_huila_productividad}
\scriptsize
\begin{tabular}{lrrr}
\toprule
Indicador & 2010 & 2024pr & Crec. anual \\
\midrule
PIB real (Billones de pesos de 2015) & 11,7 & 15,4 & 1,95\% \\
Ocupados (Millones) & 0,36 & 0,49 & 2,08\% \\
PIB por trabajador (Millones de pesos de 2015) & 32,2 & 31,6 & -0,13\% \\
Horas semanales por trabajador & 44,2 & 42,8 & -0,24\% \\
PIB por hora trabajada (Miles de pesos de 2015) & 14,0 & 14,2 & 0,11\% \\
\bottomrule
\end{tabular}
\end{table}

\subsection{Cundinamarca}
\textbf{Entre 2010 y 2024pr, el PIB por hora trabajada de Cundinamarca creció -0,14\% anual.} El PIB por trabajador creció -0,32\% anual, mientras que las horas semanales por trabajador cambiaron -0,18\% anual. Esta diferencia muestra si la productividad medida por trabajador se mueve en línea con la productividad por hora o si está afectada por cambios en la intensidad horaria.
\begin{table}[H]
\centering
\caption{Cundinamarca: PIB, ocupados, horas y productividad laboral, 2010--2024pr}
\label{tab:departamento_cundinamarca_productividad}
\scriptsize
\begin{tabular}{lrrr}
\toprule
Indicador & 2010 & 2024pr & Crec. anual \\
\midrule
PIB real (Billones de pesos de 2015) & 38,2 & 60,5 & 3,35\% \\
Ocupados (Millones) & 0,90 & 1,50 & 3,68\% \\
PIB por trabajador (Millones de pesos de 2015) & 42,3 & 40,4 & -0,32\% \\
Horas semanales por trabajador & 45,9 & 44,7 & -0,18\% \\
PIB por hora trabajada (Miles de pesos de 2015) & 17,7 & 17,4 & -0,14\% \\
\bottomrule
\end{tabular}
\end{table}

\subsection{La Guajira}
\textbf{Entre 2010 y 2024pr, el PIB por hora trabajada de La Guajira creció -0,51\% anual.} El PIB por trabajador creció -0,80\% anual, mientras que las horas semanales por trabajador cambiaron -0,29\% anual. Esta diferencia muestra si la productividad medida por trabajador se mueve en línea con la productividad por hora o si está afectada por cambios en la intensidad horaria.
\begin{table}[H]
\centering
\caption{La Guajira: PIB, ocupados, horas y productividad laboral, 2010--2024pr}
\label{tab:departamento_la_guajira_productividad}
\scriptsize
\begin{tabular}{lrrr}
\toprule
Indicador & 2010 & 2024pr & Crec. anual \\
\midrule
PIB real (Billones de pesos de 2015) & 7,5 & 8,2 & 0,62\% \\
Ocupados (Millones) & 0,27 & 0,33 & 1,43\% \\
PIB por trabajador (Millones de pesos de 2015) & 27,8 & 24,8 & -0,80\% \\
Horas semanales por trabajador & 44,7 & 43,0 & -0,29\% \\
PIB por hora trabajada (Miles de pesos de 2015) & 11,9 & 11,1 & -0,51\% \\
\bottomrule
\end{tabular}
\end{table}

{\footnotesize Nota general: PIB real en billones de pesos constantes de 2015; ocupados en millones; PIB por trabajador en millones de pesos constantes de 2015; PIB por hora en miles de pesos constantes de 2015. Fuente: cálculos propios con DANE y GEIH.}